% discussion.tex -- Discussion and Conclusion (merged to save space)

The experiments validated the core claim: Langevin dynamics on the modern Hopfield energy converts deterministic attention into a principled stochastic sampler governed by a single temperature parameter.
The temperature spectrum (Fig.~\ref{fig:experiments}a) exhibited a smooth phase transition between retrieval and generation, the continuous analogue of the classical Hopfield/Boltzmann duality~\cite{ackeleyLearningAlgorithmBoltzmann1985} lifted to continuous states and Langevin dynamics. The convergence experiment (Fig.~\ref{fig:experiments}b) confirmed that Algorithm~1 reached the correct Boltzmann target, consistent with ULA theory~\cite{durmusNonasymptoticAnalysisUnadjusted2017}. The load-ratio sweep (Fig.~\ref{fig:phase-diagram}) showed that generation quality degraded gracefully as $K/d$ increased, providing empirical evidence linking modern Hopfield capacity to generation quality. The MNIST experiment (Table~\ref{tab:mnist-baselines}, Fig.~\ref{fig:mnist-baselines}) demonstrated that these properties carried over to real images: the sampler generated novel, low-energy digit images that no static baseline could approach, including a VAE trained on the same $K{=}100$ patterns ($\mathcal{N}{=}0.214$, $\bar{\mathcal{D}}{=}0.441$, the strongest non-Langevin result): at $\beta{=}200$, SA led the VAE by $2.6{\times}$ on novelty and $2.0{\times}$ on diversity. SA matched the Metropolis-corrected MALA baseline, confirming ULA's discretization bias was negligible at $\alpha{=}0.01$. Appendices~\ref{app:finance} and~\ref{app:simpsons} extended these results to S\&P~500 log-returns ($d{=}424$) and face images ($d{=}4{,}096$), confirming the SNR rule and the novelty--fidelity trade-off across a $5.2{\times}$ dimension range.

In summary, stochastic attention reuses the same primitive operations as a standard attention head and requires no learned score network, no training loop, and no contrastive objective. The analytic structure of the Hopfield energy (smoothness, dissipativity, and an exact closed-form gradient) provides convergence guarantees that generic energy-based samplers cannot offer without additional assumptions. Across all four experimental domains ($d\in\{64,\,784,\,424,\,4{,}096\}$), the transition from structured retrieval to diffuse generation consistently occurs near $\mathrm{SNR}=\sqrt{\alpha\beta/2d}\approx 0.02$--$0.03$, giving a dimension-independent rule for targeting $\beta$: set $\beta=\mathrm{SNR}_{\mathrm{target}}^{2}\cdot 2d/\alpha$.

\paragraph{Limitations.}
The ULA discretization introduces bias scaling with $\alpha$; our MALA comparison (Table~\ref{tab:mnist-baselines}) shows a 99.2\% acceptance rate at $\alpha{=}0.01$, confirming negligible bias at this step size. At larger $\alpha$, MALA (Algorithm~\ref{alg:mala}, Appendix~\ref{app:mala-algorithm}) is the preferred variant. Mixing time grows with both $\beta$ and $K/d$, and quantitative scaling laws remain to be established. Each step costs $\mathcal{O}(NK)$ (two matrix--vector products plus a softmax), identical to a single attention head; for $N{=}784$, $K{=}100$ the full 30-chain experiment completes in under 10 minutes on a laptop CPU. The method also requires $\mathbf{X}$ to be known and fixed at sampling time; streaming or latent memory settings would require extension. We consider only single-head attention over a flat memory; extension to multi-head attention, hierarchical memories, or full encoder-decoder transformers is future work.
